%============================================================
% Snow and Voting in Austria
%============================================================
 
\documentclass{article}
 
% ── Page layout ──────────────────────────────────────────────
\usepackage[margin = 2.54cm]{geometry}
 
% ── Encoding ─────────────────────────────────────────────────
\usepackage[utf8]{inputenc}
 
% ── Tables & figures & References ─────────────────────────────────────────
\usepackage{graphicx, dcolumn}
\usepackage{booktabs}
\usepackage{tabularray}
\usepackage{float}
\usepackage{siunitx}
\usepackage[style=apa,backend=biber]{biblatex}
\addbibresource{literature.bib}
% ── Line spacing ─────────────────────────────────────────────
\usepackage{setspace}
\doublespacing
\usepackage{xurl} 
\usepackage{hanging}
\usepackage{indentfirst}
 
% ── Fonts ────────────────────────────────────────────────────
\usepackage[T1]{fontenc}
\usepackage{times}
 
% ── Paragraph indent ─────────────────────────────────────────
\setlength{\parindent}{0.5in}
 
% ── Section heading formatting ───────────────────────────────
\usepackage{titlesec}
\titleformat{\section}[block]
  {\normalfont\normalsize\bfseries}{{\Roman{section}.\quad}}{0em}{}
\titleformat{\subsection}[block]
  {\normalfont\normalsize\itshape}{}{0em}{\hspace{0.5in}}
\titleformat{\subsubsection}[block]
  {\normalfont\normalsize\bfseries\itshape}{}{0em}{\hspace{0.5in}}
\titlespacing*{\section}{0pt}{12pt}{6pt}
\titlespacing*{\subsection}{0pt}{10pt}{4pt}
\titlespacing*{\subsubsection}{0pt}{10pt}{4pt}
 
% ── Captions (APA 7) ─────────────────────────────────────────
\usepackage{caption}
\captionsetup{
  labelfont=bf,
  labelsep=newline,
  justification=raggedright,
  singlelinecheck=false,
  font=normalsize
}
 
% ── Hyperlinks ───────────────────────────────────────────────
\usepackage[hidelinks]{hyperref}
 
% ── Math ─────────────────────────────────────────────────────
\usepackage{amsmath}
 
% ── Page numbers bottom-right ────────────────────────────────
\usepackage{fancyhdr}
\pagestyle{fancy}
\fancyhf{}
\rfoot{\thepage}
\renewcommand{\headrulewidth}{0pt}
 
%============================================================
%  DOCUMENT
%============================================================
\begin{document}

% ── TITLE PAGE ───────────────────────────────────────────────
\begin{titlepage}
  \centering
  \vspace*{3cm}

  {\large\bfseries The Politics of Snow:\\ 
  Climate Change and Economic Grievances in Austrian Municipalities}

  \vspace{4cm}

  {\normalsize\ Franziska Rogg}\\
  {\normalsize\ 100580049}\\
  {\normalsize Master in Social Sciences}\\
  {\normalsize Political Economy}\\
  {\normalsize I\c{s}ik Özel, Dulce Manzano, Sebastian Lavezzolo}\\
  {\normalsize 19/06/2026}

\end{titlepage}
% ── ABSTRACT──────────────────────────────────────────
\section*{Abstract}
\noindent 
This paper examines to what extent snow scarcity in municipalities shapes electoral outcomes of general elections. While previous literature largely focuses on climate-induced natural disasters and their political consequences, I argue that continously changing snow scarcity can shape voting behaviours especially in lower-altitude tourism-dependent municipalities. I hypothesize the effect to operate through two different channels: economic grievance and environmental salience, which are reflected in either heightened far-right vote shares or increased Green party vote shares. Using municipal-level data on average seasonal snow depth in Austrian municipalities, I exploit general election outcomes of Austria's national elections in 2024, which allows to obtain behavioural data on political preferences in the context of exposure to climate change. The findings show evidence of a modest effect for increased far-right vote shares when exposed to declining snow depth, indicating mechanisms of far-right backlash due to distributive conflicts. The results show no sizeable effect of snow scarcity on Green Party support, failing to provide evidence for heightened environmental salience in contexts of high vulnerability to climate change. 
The study contributes to the literature exploring impacts of climate change and vulnerability on electoral outcomes. 
\vspace{2em}

% ── I. INTRODUCTION ──────────────────────────────────────────
\section{Introduction}
Rising temperatures, longer summers and shorter winters. As climate reports have long predicted, seasonal snowfall patterns, snow availability, and precipitation trends have proved continuous shifts throughout the Northern Hemisphere \parencite{li2025, Olefs2020, oecd2007}. Climate change therefore concentrates snowfall in higher-altitude regions as well as Western parts of the Alpine region, disproportionately burdening lower-altitude regions with snow scarcity, shorter winter seasons and higher temperatures \parencite{li2025, Steiger2013}. While compensatory measures such as snowmaking offer partial relief, their effectiveness is limited by rising temperatures that reduce suitable snowmaking days, and their implementation requires heavy investments and high operating costs \parencite{Steiger2013, copernicus2018}. The consequences of anthropogenic climate change extend beyond the natural environment and the availability of resources, involving drastic economic implications for industries built on it, such as winter tourism in Alpine economies \parencite{lealfilho2024}.
Winter tourism, long a crucial economic pillar of Alpine regions, faces severe resource pressure from declining availability of natural snow in lower altitude regions, while other, higher-altitude regions might profit in the short-term from shifting snowfall patterns and structural demand changes \parencite{oecd2007, lealfilho2024}. Austria is a central player in the European but also worldwide tourism sector, which is reflected in the domestic relevance of tourism income, making up 6.3\% of Austria's gross domestic product in 2024 \parencite{fritz2025}. With approximately 440 ski resorts forming the backbone of its winter tourism industry, more than 70\% are located at lower elevations below 1500 meters, placing the majority of Austria's winter tourism destination at high exposure to the effects of climate change on natural snow availability. 

Yet the consequences of declining snow availability extend beyond the immediate natural and economic consequences, shaping the material conditions of entire municipalities through lost tourism demand and financial burdens through compensation costs. The conventional account of economic voting shows that voters reward or punish the incumbent for the state of the economy \parencite{LewisBeck2000, Duch2006}, suggesting that municipalities, bearing a disproportionate share of climate-induced costs, may translate their economic grievance into political preferences at the ballot box. Recent findings suggest that personal experiences with climate change shape political attitudes towards environmental policies and support, as well as general election outcomes \parencite{egan2012, baccini2021, bechtel2011, Pink2023}. Furthermore, natural disasters have shown that voters hold the incumbent accountable in general elections for disaster relief measures as well as heightened environmental salience in electoral outcomes \parencite{bechtel2011, Pink2023}. While this literature largely relies on findings concerning destructive natural disasters induced by climate change, it remains underexplored how slow-moving, continuous consequences of climate change, such as declining resource availability, shape political attitudes. 
Whether such gradual climate stress channels grievances toward far-right parties demanding economic relief for the arising distributive conflict \parencite{Rho2017, Halikiopoulou_Vrakopoulos_Arndt_2026, Habersack2024, Margalit2019}, or instead heightens environmental salience and drives support for green parties \parencite{baccini2021, Pink2023} still remains contested in the literature and requires further exploration. 
Expanding on these streams of literature, this study explores the impact of seasonal snow availability in municipalities on electoral outcomes in Austria's national elections in September 2024. To what extent, if at all, does the the seasonal snow depth in municipalities affect party vote shares in general elections? Does snow scarcity induce far-right backlash and how does it shape support for pro-environmental parties? 

This study contributes to ongoing debates about climate-induced political consequences, environmental support as well as mechanisms of economic voting through material grievances and distributive conflicts. Although \textcite{Mildenberger2017} finds no trade-off between economic conditions and environmental preferences, this study sheds light on the salience of climate-induced economic grievances and far-right backlash as opposed to environmental concern in general elections, given the context of high exposure and vulnerability to climate change. 

The remainder of this article is structured as follows. In the next section, I develop the theoretical framework and argument guiding the analysis. Sections III and IV introduce the data and empirical strategy as applied in the following analysis. In Section V I  present the results, which are further discussed in Section VI before Section VII condludes. 




\newpage
% ── II. Theoretical Framework ──────────────────────────────────
\section{Theoretical Framework}
The paper builds on rich streams of literature in the field of political economy that evolved around climate change and political attitudes, environmental policy support as well as economic voting, distributive conflict, and populist backlash.  Below I provide a brief overview of each literature body relevant to this study as well as my theoretical argument. 

\subsection{Climate Change \& Political Attitudes}
This body in the literature largely builds on findings concerning political attitudes shaped by personal experiences related to climate change and natural disasters. Several studies find evidence of a positive link between personal experience and environmental attitudes. \textcite{egan2012}  and \textcite{Li2011} find that unusual anomalies of local temperatures strengthen personal beliefs about global warming in the short-run, especially in contexts of low education and weak party affiliations.  Building on this, recent papers found that proximity to natural disasters positively affects beliefs about global warming and support for climate policies \parencite{baccini2021, Konisky2015}. Research in this stream shows that personal exposure to climate-related events or rising temperatures temporarily increases both climate change concern and support for environmental policies in the immediate aftermath. 

Further research links natural disasters to political behaviour. The guiding framework  for this work is \textit{retrospective voting} which posits that rational voters hold incumbents responsible for relief performance and reconstruction efforts in the aftermath of natural disasters, in addition to positive economic performance in times without disaster \parencite{HEALY2009, Achen2016-zp}. Evidence from the municipal-level shows that voters reward the incumbent with re-election for positive reconstruction and relief efforts immediately after the disaster, showing signs of institutional resilience and advantage of incumbents when it comes to effective measures to alleviate economic consequences of the disaster  \parencite{Masiero2021, TaheriHosseinkhani2026}. While \textcite{bechtel2011} find further evidence of \textit{voter gratitude}, where voters reward the incumbent with re-election, \textcite{Pink2023} show that voters rationally evaluate disaster management and punish malperformance of candidates while also increasing support for the Green Party. The increase in sensitivity of environmental concern is in line with the findings of \textcite{baccini2021} and \textcite{Masiero2021}. Yet, \textcite{Blankenship2020} establish in their study, that retrospective voters punish only opposition leaders, who are in government but not in line with the government leader and therefore only have limited power to provide effective relief measures. Similarly, \textcite{Gasper2011}  argue that attentive voters are able to attribute responsibility to the respective politician and therefore punish in a directed matter. As these findings show electoral incentives for disaster responsiveness by local governments \parencite{Livert2025}, \textcite{HEALY2009} draws certain political economy implications from these observed patterns of short-lived reward for relief efforts. When voters reward disaster relief, politicians are incentivized to increase investments in relief aid measures and less so in disaster preparedness. While this is electorally efficient from a policy perspective, it is less effective from an economic perspective, especially in times with rising frequency of natural disasters and progressing climate change \parencite{Seneviratne2023}. 

This stream in the literature effectively shows how personal exposure to climate change and natural disasters can shape environmental attitudes, as well as political outcomes for the incumbent. However, the study of personal experience mostly draws on survey data to evaluate attidudes, which comes short of providing insights into how these attitudes translate into effective electoral outcomes. Although the studies of natural disasters and political outcomes draw on electoral outcomes, the results are limited to disasters as exogenous shocks, missing the effects of climate change over time, such as the continuous decrease in snow availability in resource-dependent regions. 
The theoretical foundation of this literature is largely rooted in retrospective and economic voting, connecting to the role of domestic distributive conflicts and populist backlash.  


\subsection{Economic Voting \& Backlash}

Personal experience is central to how voters form political preferences, a crucial mechanism found in the literature on economic voting \parencite{Duch2006}. Especially for less informed and engaged voters, experience represents a low-effort and immediate source of information that influences voting behaviour. 
Scholars have long established that economic conditions shape political outcomes, as the economic voter holds the government responsible for its fiscal policies and the state of the economy \parencite{LewisBeck2000, LEWISBECK2011288}. \textcite{Fossati2014} finds evidence for an increase in economic prevalence, for voters who are most vulnerable to unemployment. Furthermore, \textcite{Anderson2000} specifies that voters are even more capable of expressing economic discontent when blame can be attributed transparently. 
Although voters evaluate the economy when forming political opinions, they do not take the same position on how to improve contemporaneous economic conditions \parencite{LEWISBECK2011288}. As \textcite{Habersack2024} find that the perception of one's personal or the societal economic position is associated with a \textit{greater distance from the government}. For voters who are worse off economically, this observation drives increasing populist attitudes in Western democracies.
An important implication that most of the above share, is that economic preferences are short-lived, therefore, mostly present in immediate elections. 

The mechanism of economic discontent driving populist voting is far from novel, but has been largely discussed in the literature concerning globalization, redistribution and populist backlash.
\textcite{rodrik2018} argues that as globalization progressed, economic anxiety and discontent created pathways for right-wing parties in Europe that opposed the liberal political mainstream and created a simple narrative for who to blame. \textcite{Rho2017}
find the prevalence of a mechanism of economic ignorance that causes a cleavage between liberal preferences and material self-interest. By showing how beliefs affect economic preferences, \textcite{Rho2017} argue that voter preferences might not always reflect their best economic interest, but may rather be shaped by economic anxiety. Although the role of economic insecurity in shaping populist backlash is well established in the literature, \textcite{Margalit2019} argues for a more holistic understanding of the economic and non-economic influences, such as culture and immigration, that shape the populist vote. Further research suggests that the causal relationship between economic grievances and the rise of populism remains contested \parencite{Mudde2018}. \textcite{Habersack2024} concludes that while economic hardship plays a significant role in the populist vote, political factors and cultural channels are further important contributors. 


\subsection{Vote Choice}
A growing body of literature has established various links between environmental concerns, climate change policy support, and backlash. While \textcite{patterson2023} established a framework for the backlash against climate politics, further research established a conditional link between economic circumstances and environmental support, as well as spatial cleavages in environmental protection \parencite{Mildenberger2017, Halikiopoulou_Vrakopoulos_Arndt_2026, Arndt2023}. \textcite{Halikiopoulou_Vrakopoulos_Arndt_2026} argue that rural voters tend to electorally resist green policies as they perceive to be burdened with the costs of such policies. Conversely, green voters tend to concentrate in rather urban areas and are less likely to be burdened with disproportionately high transition costs. Furthermore, scholars have long argued that in times of economic downturn, environmental preferences weaken, as individuals evaluate short-term fiscal needs of higher priority \parencite{Shum2012}. Yet, \textcite{Mildenberger2017} find no such effect at the individual level, suggesting that rather sociotropic economic changes condition aggregated environmental beliefs as opposed to individual economic downturns. 
In the Austrian context, \textcite{kritzinger2013} argue that electoral choices are shifting away from traditional perceptions of class-based loyalties towards issue-based decisions that are shaped by party framing and campaign narratives. As research showed, European electoral outcomes are shaped by territorial cleavages between rurally concentrated far-right voting, whereas green voters tend to concentrate in urban areas \parencite{Arndt2023, DeVries2018, Halikiopoulou_Vrakopoulos_Arndt_2026}. 
However, it remains underexplored how climate vulnerability influences these determinants of vote choice given the salience of the tourism industry in Austria. 
The far-right \textit{Freiheitliche Partei Österreich} (FPÖ) achieved the highest party vote share in the national elections 2024, by proposing a strict \textit{'Austria First'} stance in the economy, social welfare, immigration, as well as foreign policy. Although the party acknowledges climate change, it opposes costly green transitions by framing it as an elite-driven policy effort that disproportionately burdens economically vulnerable voters \parencite{fpoe2024}. 
Conversely, the Green Party represents the main promoter of environmental and climate policy, centering their electoral program around the \textit{'Green and Social Deal'}. The party pledges for renewable energy transitions as well as increased state spending to enforce social welfare and the creation of green jobs \parencite{gruene2024}.

The discussions in the literature concerning economic voting, far-right backlash, as well as environmental support are well established. However, the effect of climate change on the natural resource base that shapes the profitability of certain industries, as well as its electoral consequences yet remain to be explored. 


\subsection{Theoretical Argument}
As shown above, the previous literature is well established in the strands of climate change and political attitudes, economic voting, populist backlash and determinants of vote choice. This study is positioned at the intersection of these literature streams, as the connection of how changing resource-availability, conditioned by climate change, shapes electoral outcomes through mechanisms of economic voting and environmental concern remains underexplored.
In the following, I will argue for two possible channels through which snow availability shapes voting behaviour.
Firstly, as climate change proceeds to impact snowfall patterns across the Alpine regions, differences in the seasonal snow cover between municipalities will be reflected in local voting behavior in Austria's national elections 2024. Declining snow depth is reflected in shorter winter seasons, especially in lower-altitude municipalities. As rural municipalities, often specializing in winter tourism, tend to suffer from reduced snow availability, arising economic anxiety and grievance are reflected in spatially concentrated far-right vote shares on the municipal level. Declining snow availability directly threatens the profitability of tourism as an important income base, by increasing compensation costs or by reducing tourism demand in affected municipalities. These costs are not distributed evenly, giving rise to a climate-induced distributive conflict. Higher-altitude municipalities can offset revenue losses through artificial snow production and by attracting tourism demand redirected from snow-scarce destinations. Lower-altitude municipalities with high tourism dependence and limited fiscal capacity, by contrast, bear a disproportionate share of the adjustment burden — concentrating economic grievances where snow scarcity is high. Snow scarcity is reflected in shorter and less reliable winter seasons as well as high dependence on snowmaking. While this increases operation costs disproportionately for lower-altitude municipalities, its success is further limited by low temperatures and water availability \parencite{Steiger2013}. Economic grievances rooted in this unequal distribution of costs are therefore spatially concentrated in vulnerable municipalities with less snow availability. Rural and vulnerable voters, therefore, tend to prioritize short-term fiscal needs rather than costly green policies and energy transition that pose an additional economic burden. Faced with economic precarity, voters oppose costly environmental campaigns and turn toward the far-right as a political outlet, aiming for short-term economic relief. 
Conversely, reduced snow availability may increase the electoral salience of environmental concerns among voters immediately vulnerable to climate change. As voters are immediately exposed to consequences of climate change, personal attitudes toward global warming shift towards increased environmental salience, driving support toward the Green Party as the major pro-climate party in the political system. 
These effects are expected to differ further between higher and lower altitude municipalities as rising temperatures and precipitation patterns increase the snowline and concentrate snowfall in higher elevations \parencite{li2025}. Furthermore, municipalities of lower altitude are less capable to engage in compensation efforts such as artificial snowmaking or snowfarming, as these require sustained cold temperatures and expensive investments in specific technologies. 
Based on the main opposing mechanisms of economic grievance or environmental salience, this study aims to explore the effect of seasonal natural snow availability on party vote shares in the 2024 Austrian elections. The following hypotheses are therefore put forward: 
\vspace{1em}

\noindent $H_1$: Municipalities with less snow depth are associated with higher vote shares of the FPÖ.

\vspace{0.5em}
\noindent $H_2$: Municipalities with less snow depth are associated with higher vote shares of the Green Party.

\vspace{0.5em}
\noindent $H_3$: Municipalities with negative deviations in snow depth from historical averages are associated with higher vote shares of the FPÖ.

\vspace{0.5em}
\noindent $H_4$: Municipalities with negative deviations in snow depth from historical averages are associated with higher vote shares of the Green Party.

\vspace{0.5em}
\noindent $H_5$: The negative association between snow depth and far-right vote shares is expected to weaken in municipalities of higher altitude.



\vspace{2em}




% ── III. DATA ──────────────────────────────────
\section{Data}
% Template for adding a subsection if needed:
% \subsection{Subsection Title}

To test the above-mentioned hypotheses, the case of Austria will be leveraged for the following reasons. First, the Geosphere Snowgrid dataset covers the daily snow depth of Austria in a 1 x 1 raster, providing good and reliable municipal-level data of snow depth across the country. This study used entries ranging from 1984 to 2010 as well as the years of 2023 and 2024. Furthermore, as Figure 1 shows, the winter season of 2023/2024 followed common trends in snowfall of recent years, where snowfall patterns have produced consistently higher snowfall in the Western Alps and in higher elevation than in lower provinces along the Eastern Alps \parencite{avalanchereport2024, Olefs2020, li2025}. These meteorological developments result in higher levels of seasonal snow depth in Western municipalities than in the Eastern Alps, as well as contiuous snow scarcity in lower elevations \parencite{li2025}. This provides us with differences in the extent of the seasonal snow cover across Austrian municipalities within the same institutional and electoral but also economic context, allowing for a comparison of voting behaviour under varying snow conditions. Figure 2 further illustrates this development of snow depths per winter season by showing the extent of deviations from the historical average in snow depth per municipality. 
Second, in September 2024 Austria held its National Elections, providing aggregated, behavioural data on general voter attitudes across all municipalities of Austria at the same time. This allows to study the effect of snowfall on voting behaviour across municipalities and also geographic regions with varying extents of snow depth. However, the explanatory power of national election results in regards of climate vulnerability remains limited, given that voters care about a variety of policies and do not exclusively vote based on developments in snowfall. 
Third, Austria ranks among the leading countries of Europe in terms of tourist arrivals, making the tourism industry an important source of income, especially in the Alpine provinces of the country \parencite{fritz2025}. This allows for the assumption that especially residents of tourism-intense regions are economically dependent on certain levels of snow depth per winter season but also on persistent tourism influx in the short run and might reflect this in their voting behaviour. 

The main outcome variable is voting behaviour, measured as party vote shares in the national elections in 2024 that are aggregated at the municipal-level as provided by the Austrian Ministry of the Interior \parencite{bmi2024}. In line with the hypothesized mechanisms, the analysis focuses on the respective vote shares of the two parties: the Green Party as well as the far-right FPÖ, which emerged as the strongest party in the 20240 elections  with 28.9\% of the vote. Both are relevant parties in the Austrian party system that draw on fundamentally different voter bases - the Greens predominantly rely on liberal-urban voters with strong environmental preferences, while the FPÖ draws on rural conservatives but also economically vulnerable workers \parencite{kritzinger2013}. The workforce of the tourism industry is mainly shaped by a duality of vocationally skilled, seasonal workers, such as waitstaff, ski instructors and lift operators, opposed to locally-rooted executives, landowners and hoteliers.  
Although the FPÖ acknowledges climate change, the party presents itself skeptical of its anthropogenic nature and cost-bearing energy transition policies, rather campaigning for policies enhancing industrial and short-term economic growth, among others \parencite{fpoe2024, ennserjedenastik2016}.
While this approach does not allow to isolate explicit grievances from reduced snow depth in the previous winter season from broader voting behaviour, ideology, culture or other economic preferences, it provides a first approximation of the extent to which this exposure to external forces, such as snowfall, shapes electoral behaviour in the context of national elections. 

The key independent variable, Figure 1, is the mean depth of snow in each Austrian municipality in the winter season of 2023/2024, as provided by the Geosphere Snowgrid Raster \parencite{Olefs2020}. The dataset captures the daily extent of the snow layer in meters in a 1x1 km raster of Austria. To make it comparable to vote shares per municipality the snow depth has been aggregated to municipal-averages within the time period of a 'normal' winter season, lasting from November to April the following year. This time frame is based on usual operating times of tourism infrastructure, such as cable cars, but also hotels and other tourism corporations. This operationalization of snow availability seems adequate in the context of this study as the main outcome is general voting behaviour in national elections that take place in the same year but not in the direct aftermath of the winter season of interest. Therefore, it seems reasonable that individuals, aggregated at the municipal-level, draw on general perceptions of snow availability in the last winter season, reflected in a seasonal average, rather than averages of specific months or specific days with maximum snow depth. However, for robustness purposes, the model specification employing the maximum snow depth of the winter season 23/24 is included in the Appendix. 
To further capture climate change as a shock rather than a gradual trend, I additionally employ a snow depth anomaly measure, calculated as the deviation of the 2023/24 mean snow depth from the historical long-run municipal average, such that negative values indicate below-average snow conditions, in meters, and positive values indicate above-average snow depth (Figure 2). The historical average is based on the municipal-level aggregation of mean snow depths per winter season within the time period between 1984 - 2010, following examples of snow reports and other literature \parentcite{Olefs2020}. 

Apart from the main variables, other indicators have been collected to control for various confounding factors that vary between municipalities, such as: population size, share of the population over 65, size of the foreign population, unemployment rate, share of tertiary education and turnout in the 2024 elections \parencite{statistikaustria2025}. These were selected based on previous literature, providing insights into factors that largely drive vote share patterns \parencite{hooghe2026, ennserjedenastik2016}. 
Furthermore, to control for another geographical factor that could influence differences in snow depth per municipality, the altitude for each municipality is controlled for.
Summary statistics for the main variables are included in the Appendix (Appendix A).

\vspace{2em}

% ── IV. Methodology ────────────────────────────────
\section{Methodology}

To identify the effect of differences in snow depth between municipalities, the study aims to compare municipalities with higher depth of snow to ones with lower snow depth, approximated by specifying a cross-sectional OLS regression with district fixed effects. This model specification seems appropriate as both, the election and the winter season of interest, are within the same year, leaving no time variation to exploit. Given the context, district fixed effects are preferred over regions or states, as they provide a more thorough control, exploiting variation between municipalities within a similar economic, administrative and cultural context of the respective district. As districts are smaller and more homogeneous units than regions or states, the fixed effect absorbs effectively all unobserved time-invariant characteristics such as income levels, economic structures, grievances from Covid-19 pandemic, location or local voting traditions that could confound the relationship between snow conditions and voting outcomes. Snow depth as a predictor can be treated as plausibly exogenous to voting behaviour as it is mainly determined by meteorological and climatic conditions that are largely independent of direct political processes. 
Municipal-level controls as described above are included to account for characteristics that might differ between municipalities of the same district, potentially influencing the relationship of interest. The effect of snow depth on vote shares is estimated the following way: 
\vspace{1em}
\begin{equation*}
V_{ij} = \beta_1 \text{Snow}_{ij} + \beta_2 X_{ij} + \alpha_j + \epsilon_{ij}
\end{equation*}

\noindent where $V_{ij}$ is the vote share of the respective party in municipality $i$ in district $j$, $\text{Snow}_{ij}$ represents the snow depth coefficient, $X_{ij}$ is a vector containing the employed control variables, $\alpha_j$ accounts for district fixed effects, and $\epsilon_{ij}$ is the residual error term. 
Therefore, the main model aims to estimate the effect of within-district variations in snow depth on party vote shares across municipalities. 
The outcome variables, vote shares of the FPÖ and the Green Party in the national elections 2024, are estimated in separate models, as well as the two main independent variables, the mean snow depth 23/24 and its deviation from historical averages.
To test whether the effect of snow depth on voting behaviour varies across altitude, an interaction term between mean snow depth and altitude of the municipality is additionally estimated. This spatial heterogeneity analysis allows for an assessment of whether the tested relationship differs between municipalities of high or low altitude. 
Assessing the sensitivity of the main results, the mean snow depth is replaced by the seasonal maximum snow depth in 23/24 as an alternative operationalization of winter snow conditions. The results of this robustness check are included in Appendix B in order to verify that the results are not driven by the aggregation of seasonal snow depth per municipality. 
The unit of analysis is the municipality, providing a sample of 2,078 observations within 79 districts. Districts consisting of only a single municipality are excluded from the analysis, as fixed effects cannot be identified without within-district variation. Standard errors are clustered at the district level to correct for correlation of the residuals. 

\begin{figure}[H]
  \centering
  \includegraphics[width=0.85\textwidth]{plots/snow_mean_map.png}
  \caption{\textit{Mean Snow Depth Across Municipalities (2023/24)}}
  \label{fig:mean_map}
\end{figure}

\begin{figure}[H]
  \centering
  \includegraphics[width=0.85\textwidth]{plots/snow_anomaly_map.png}
  \caption{\textit{Snow Depth Anomaly Across Municipalities (2023/24)}}
  \label{fig:anomaly_map}
\end{figure}
\vspace{2em}

%\begin{figure}[H]
 % \centering
  %\includegraphics[width=0.85\textwidth]{plots/fpö_share_map.png}
  %\caption{\textit{FPÖ Vote Share Across Municipalities}}
  %\label{fig:fpoe_map}
%\end{figure}


\begin{figure}[H]
  \centering
  \includegraphics[width=0.85\textwidth]{plots/boxplot_fpö_snow.png}
  \caption{\textit{FPÖ Vote Share by Snow Depth Group}}
  \label{fig:boxplot}
\end{figure}
\vspace{2em}

%\begin{figure}[H]
%  \centering
 % \includegraphics[width=0.5\textwidth]{plots/vote_shares_bar.png}
%  \caption{\textit{Mean Vote Shares: FPÖ and Greens}}
 % \label{fig:barplot}
%\end{figure}


% ── V. RESULTS ──────────────────────────────────────────────
\section{Results}
 \subsection{Identifying the effect of snow depth on party vote shares}
 
As Figure 3 descriptively shows, the median far-right vote shares of municipalities within the bottom quartile of snow depth are marginally higher in comparison to the top quartile of snow depth. Furthermore, the inter-quartile range is narrow and comparable in both groups, suggesting only limited variation in vote shares across municipalities within this range, while outliers are present at both ends. To assess whether this pattern holds statistically, the specified models estimate the association between snow depth and party vote shares. 
For all model specifications, OLS regressions without district fixed effects were estimated, however, the coefficients consistently shrink upon the inclusion of fixed effects. This indicates omitted variable bias in the pooled OLS estimate and supports the fixed effects model as the preferred model for description and interpretation. Additionally, the inclusion of district fixed effects substantially improves the model fit, with R² increasing from 0.378 to 0.597, underscoring the importance of controlling for unobserved time-invariant district-level characteristics.

As Table 1 shows, the coefficients on the mean seasonal snow depth are negative for vote shares of both parties. The association between average seasonal snow depth and FP Ö vote shares is moderately significant at the 5\% \ level, translating into a 3.25 percentage point decrease in far-right vote shares on average for each one-meter increase in snow depth within the district and keeping all other factors constant. Conversely, these results show a modest increase in far-right vote shares as the average snow depth decreases. This negative association is consistent with H1, suggesting that municipalities with less natural snow availability throughout the winter season tend to show higher far-right support.
The effect on Green vote shares fails to reach statistical significance at conventional significance levels. Although the direction of the effect is as expected, the results, however, are inconsistent with H2 due to statistical insignificance.

Evaluating the shock of changes in snow depth compared to historical values per municipality, Table 2 shows that the coefficients on anomaly are negative for both party specifications. Although neither of the coefficients reaches conventional statistical significance, the direction is in line with the results based on the seasonal mean snow depth, suggesting the null finding reflects limited statistical power based on the presented data rather than the absolute absence of an effect. 

\vspace{2em}
\begin{table}[H]
  \caption{\textit{Effect of Mean Snow Depth on Vote Shares}}
  \label{tab:mean}
  \begin{table}
\centering
\begin{talltblr}[         %% tabularray outer open
entry=none,label=none,
note{}={+ p \num{< 0.1}, * p \num{< 0.05}, ** p \num{< 0.01}, *** p \num{< 0.001}},
]                     %% tabularray outer close
{                     %% tabularray inner open
colspec={Q[]Q[]Q[]Q[]Q[]},
hline{2}={1-5}{solid, black, 0.05em},
hline{20}={1-5}{solid, black, 0.05em},
hline{1}={1-5}{solid, black, 0.08em},
hline{23}={1-5}{solid, black, 0.08em},
column{2-5}={}{halign=c},
column{1}={}{halign=l},
}                     %% tabularray inner close
& OLS FPÖ & FE FPÖ & OLS Greens & FE Greens \\
(Intercept) & \num{59.120}*** &  & \num{9.695}*** &  \\
& (\num{4.274}) &  & (\num{1.284}) &  \\
mean\_2324 & \num{-9.610}*** & \num{-3.249}* & \num{0.674}* & \num{-0.267} \\
& (\num{0.992}) & (\num{1.383}) & (\num{0.276}) & (\num{0.378}) \\
altitude & \num{0.005}*** & \num{0.002} & \num{-0.001}*** & \num{-0.001}* \\
& (\num{0.001}) & (\num{0.001}) & (\num{0.000}) & (\num{0.000}) \\
turnout & \num{-0.246}*** & \num{-0.254}*** & \num{-0.084}*** & \num{-0.038}** \\
& (\num{0.040}) & (\num{0.044}) & (\num{0.012}) & (\num{0.014}) \\
log\_pop & \num{0.815}*** & \num{0.425}+ & \num{0.186}** & \num{0.219}** \\
& (\num{0.199}) & (\num{0.243}) & (\num{0.060}) & (\num{0.076}) \\
pop\_over65 & \num{-0.252}*** & \num{-0.467}*** & \num{-0.101}*** & \num{0.029}* \\
& (\num{0.038}) & (\num{0.051}) & (\num{0.011}) & (\num{0.013}) \\
pop\_foreign & \num{-0.166}*** & \num{-0.133}** & \num{0.020}+ & \num{0.007} \\
& (\num{0.030}) & (\num{0.040}) & (\num{0.010}) & (\num{0.010}) \\
unemp\_rate & \num{-0.062} & \num{-0.104} & \num{-0.200}*** & \num{-0.016} \\
& (\num{0.083}) & (\num{0.111}) & (\num{0.024}) & (\num{0.025}) \\
share\_tertiary\_edu & \num{-0.711}*** & \num{-0.656}*** & \num{0.371}*** & \num{0.398}*** \\
& (\num{0.025}) & (\num{0.045}) & (\num{0.011}) & (\num{0.014}) \\
Num.Obs. & \num{2115} & \num{2078} & \num{2115} & \num{2078} \\
R2 & \num{0.378} & \num{0.597} & \num{0.619} & \num{0.767} \\
R2 Adj. & \num{0.376} & \num{0.580} & \num{0.618} & \num{0.757} \\
\end{talltblr}
\end{table}

\end{table}
\vspace{2em}


\begin{table}[H]
  \caption{\textit{Effect of Snow Anomaly on Vote Shares}}
  \label{tab:anomaly}
  \begin{table}
\centering
\begin{talltblr}[         %% tabularray outer open
entry=none,label=none,
note{}={+ p \num{< 0.1}, * p \num{< 0.05}, ** p \num{< 0.01}, *** p \num{< 0.001}},
]                     %% tabularray outer close
{                     %% tabularray inner open
colspec={Q[]Q[]Q[]Q[]Q[]},
hline{2}={1-5}{solid, black, 0.05em},
hline{20}={1-5}{solid, black, 0.05em},
hline{1}={1-5}{solid, black, 0.08em},
hline{23}={1-5}{solid, black, 0.08em},
column{2-5}={}{halign=c},
column{1}={}{halign=l},
}                     %% tabularray inner close
& OLS FPÖ & FE FPÖ & OLS Greens & FE Greens \\
(Intercept) & \num{61.083}*** &  & \num{9.724}*** &  \\
& (\num{4.350}) &  & (\num{1.306}) &  \\
anomaly & \num{-11.115}*** & \num{-1.048} & \num{-0.783} & \num{-0.647} \\
& (\num{2.152}) & (\num{2.878}) & (\num{0.549}) & (\num{0.910}) \\
altitude & \num{0.001} & \num{-0.000} & \num{-0.000}+ & \num{-0.001}** \\
& (\num{0.000}) & (\num{0.001}) & (\num{0.000}) & (\num{0.000}) \\
turnout & \num{-0.267}*** & \num{-0.254}*** & \num{-0.083}*** & \num{-0.038}** \\
& (\num{0.041}) & (\num{0.045}) & (\num{0.013}) & (\num{0.014}) \\
log\_pop & \num{0.961}*** & \num{0.435}+ & \num{0.152}* & \num{0.222}** \\
& (\num{0.201}) & (\num{0.256}) & (\num{0.060}) & (\num{0.077}) \\
pop\_over65 & \num{-0.240}*** & \num{-0.460}*** & \num{-0.109}*** & \num{0.028}+ \\
& (\num{0.039}) & (\num{0.052}) & (\num{0.011}) & (\num{0.014}) \\
pop\_foreign & \num{-0.228}*** & \num{-0.135}*** & \num{0.023}* & \num{0.006} \\
& (\num{0.030}) & (\num{0.039}) & (\num{0.011}) & (\num{0.010}) \\
unemp\_rate & \num{-0.074} & \num{-0.155} & \num{-0.175}*** & \num{-0.012} \\
& (\num{0.083}) & (\num{0.114}) & (\num{0.024}) & (\num{0.024}) \\
share\_tertiary\_edu & \num{-0.699}*** & \num{-0.662}*** & \num{0.373}*** & \num{0.398}*** \\
& (\num{0.025}) & (\num{0.045}) & (\num{0.011}) & (\num{0.014}) \\
Num.Obs. & \num{2115} & \num{2078} & \num{2115} & \num{2078} \\
R2 & \num{0.348} & \num{0.594} & \num{0.619} & \num{0.767} \\
R2 Adj. & \num{0.346} & \num{0.576} & \num{0.617} & \num{0.757} \\
\end{talltblr}
\end{table}

\end{table}
\newpage


\subsection{Effect heterogeneity across altitude}
So far, the average effect of snow depth on party vote shares across municipalities, holding time-invariant characteristics and other controls constant, has been estimated. As argued above, the assumed mechanism for increased  far-right vote shares as the snow availability decreases are the economic grievances  municipalities suffer from when the lack of snow leads to decreased profitability of winter tourism. A way to validate this assumption is to see if the suspected negative effect is less strong in areas of higher altitude, as recent literature shows a geographic concentration of snowfall in regions of higher elevation, leading to less vulnerability to snow availability in municipalities of higher altitude \parencite{li2025}. To gauge the potential heterogeneity of this effect across altitude, Table 3 shows the coefficient on mean snow depth when holding altitude constant at zero. This coefficient itself would be nonsensical in the context of Austria, as the lowest municipalities are all above sea level. However, the interaction term allows to show the marginal effect an additional meter in altitude has on the effect of increasing snow depth on far-right vote share. As Table 3 shows both coefficients, mean snow depth as well as the interaction with altitude, are negligible and statistically insignificant. The absence of a statistically significant effect concludes that within the scope of this study the effect of snow depth does not vary with increasing altitude, inconsistent with H5. 

Taken together, the results present to be inconsistent with H2 - H5, while presenting modest evidence that is consistent with H1. This concludes to a modest negative relationship between the average snow depth per municipality in the winter season 23/24 and the far-right vote shares in the national elections, held in September 2024. Further insights into the spatial heterogeneity, the shock of changes in snow cover, as well as the extent to which Green Party vote shares are shaped by the levels of snow depth fail to reach statistical significance. 

\vspace{2em}

\begin{table}[H]
  \caption{\textit{Heterogeneity Across Altitude: Interaction Effects}}
  \label{tab:interaction}
  \begin{table}
\centering
\begin{talltblr}[         %% tabularray outer open
entry=none,label=none,
note{}={+ p \num{< 0.1}, * p \num{< 0.05}, ** p \num{< 0.01}, *** p \num{< 0.001}},
]                     %% tabularray outer close
{                     %% tabularray inner open
colspec={Q[]Q[]Q[]},
hline{2}={1-3}{solid, black, 0.05em},
hline{22}={1-3}{solid, black, 0.05em},
hline{1}={1-3}{solid, black, 0.08em},
hline{25}={1-3}{solid, black, 0.08em},
column{2-3}={}{halign=c},
column{1}={}{halign=l},
}                     %% tabularray inner close
& OLS FPÖ & FE FPÖ \\
(Intercept) & \num{59.221}*** &  \\
& (\num{4.279}) &  \\
mean\_2324 & \num{-10.358}*** & \num{-2.971} \\
& (\num{2.839}) & (\num{4.097}) \\
altitude & \num{0.005}*** & \num{0.002} \\
& (\num{0.001}) & (\num{0.001}) \\
turnout & \num{-0.247}*** & \num{-0.254}*** \\
& (\num{0.040}) & (\num{0.044}) \\
log\_pop & \num{0.816}*** & \num{0.426}+ \\
& (\num{0.198}) & (\num{0.238}) \\
pop\_over65 & \num{-0.250}*** & \num{-0.468}*** \\
& (\num{0.038}) & (\num{0.052}) \\
pop\_foreign & \num{-0.164}*** & \num{-0.134}** \\
& (\num{0.032}) & (\num{0.041}) \\
unemp\_rate & \num{-0.073} & \num{-0.101} \\
& (\num{0.089}) & (\num{0.119}) \\
share\_tertiary\_edu & \num{-0.712}*** & \num{-0.656}*** \\
& (\num{0.025}) & (\num{0.045}) \\
mean\_2324 × altitude & \num{0.000} & \num{-0.000} \\
& (\num{0.002}) & (\num{0.002}) \\
Num.Obs. & \num{2115} & \num{2078} \\
R2 & \num{0.378} & \num{0.597} \\
R2 Adj. & \num{0.376} & \num{0.579} \\
\end{talltblr}
\end{table}

\end{table}
\newpage

% ── VI. Discussion & Limitations ────────────────────────────────────────────
\section{Discussion of the Results}
\subsection{Snow Scarcity and Far-right Support}
I argue that seasonal snow availability per municipality affects electoral outcomes in Austria's national elections in September 2024. The conceptual framework above is based on rural municipalities being economically vulnerable to snow availability due to high dependence on winter tourism. Less seasonal snow depth then marginally affects vote choice at the municipal-level. One potential mechanism through which snow depth can affect vote shares are economic grievances that increase far-right vote shares. Another, is environmental salience that increases Green Party vote shares. 

The results support the argument that municipalities with lower seasonal averages in snow depth reflect slightly higher far-right vote shares. 
\textcite{Olefs2020} show that a thinner snow cover translates into a decreasing number of snow days per winter season. When confronted with a shorter natural winter season, municipalities are burdened with either high compensation costs to i.e increase investment in artificial snow production or snowfarming, or with the loss of profitable snow days. Therefore, as  seasonal snow depth decreases, its capacity to form a reliable foundation for high-intensity winter tourism decreases, leaving residents and workers in a situation of economic precarity. The economic voting literature established that voters hold the government responsible for one's individual but also societal economic position \parencite{Duch2006, LEWISBECK2011288, Habersack2024}. As municipalities with less seasonal snow depth show higher far-right vote shares, the analysis shows that voters are likely to punish the incumbent government for contemporaneous grievances when holding other potential sources of discontent, such as immigration and unemployment, constant. 

As the empirical analysis shows no significant effect on Green vote shares, there is no indication of changing environmental preferences or environmental salience as snow availability changes within districts. The absence of a significant effect may reflect several underlying explanations. Figure 4 shows the intensity of the Green vote shares in all municipalities, indicating that the most intense regions are concentrated in urban centers or regions that show a diversified industry landscape. This is consistent with findings on the Austrian electoral base \parencite{kritzinger2013}.
The green voter base is therefore less structurally exposed to declining snow availability and the economic precarity it entails, suggesting that the affected population, as hypothesized in Alpine tourism-dependent municipalities does not align with typical green electorate. Consequently, no change in environmental preferences is reflected in the electoral outcomes that is related to seasonal snow depth. Additionally, the impact of changes in snow availability might not be strong or visible enough to show heightened environmental preferences and support for climate-oriented policymaking. Although the snow cover decreased drastically over the last decades \parencite{Olefs2020}, the gradual and slow-moving nature of this decline may prevent voters from attributing it to a specific political cause or responsibility or translating it into a concrete pro-environmental vote choice. 

This evidence is in line with the findings of \textcite{Halikiopoulou_Vrakopoulos_Arndt_2026}, establishing a spatial cleavage in voting preferences. They find that economically vulnerable voters in rural areas tend to show the least support for parties that advertise  environmental policies as they burden disproportionately low-income earners, but rather prefer far-right parties as a form of electoral resistance. Furthermore, the significant effect on far-right vote shares and the negligible effect on green vote shares shows further evidence for the economy-environment tradeoff \parencite{Shum2012, Mildenberger2017}. The preceding electoral term in Austria was shaped by pandemic fatigue, that severely affected tourism revenues, as well as overall times of economic downturn \parencite{Praprotnik2025}. Consequently,  environmental preferences might be less salient in the short-term in comparison to economic grievances that drive far-right voting. This is further supported by \textcite{Zaller1992}, who argue that voters have multiple, often conflicting preferences on many political concerns. In times of economic downturns, environmental preferences decrease in significance as material insecurities and distributive conflicts become more salient and therefore present in electoral outcomes. While voters might still be concerned with the environment and climate change, short-term fiscal needs prevail during the vote choice while it possibly contradicts environmental beliefs. This is in line with \textcite{Rho2017}, who argue that beliefs do not always reflect voters' best economic interest, but may rather be shaped by economic anxiety. Winter tourism can only be sustained in the long-term by protecting the resource base it is built upon with drastic green policies. However, in the short-term, material insecurity incentivizes support for less environmentally sustainable economic relief measures.  

Furthermore, the negligibility and insignificance of the interaction effect of altitude does not directly conclude to the absence of an effect in general. The spatial heterogeneity assumes that lower-altitude municipalities do not compensate for the lack of natural snow by producing more technical snow and therefore direct tourism streams toward higher-altitude destinations. However, in reality, especially lower-altitude winter tourism destinations rely on technical snow production, snowfarming or other compensatory technologies to circumvent income losses through shorter natural winter seasons and less natural snow availability. Around 77\% of Austria's winter tourism destinations are concentrated in lower-altitude elevations and therefore largely dependent on snowmaking \parencite{Steiger2013}. Although the compensatory success of snowmaking is constrained by water availability, temperatures and high operating costs, these limits might not yet be reached, making winter tourism still profitable in lower-altitude destinations. Furthermore, these measures entail heavy investments in specific technology, the additional economic burden might not yet be attributable to political decisions or climate change. Another reason for the lack of spatial heterogeneity could be that changes in snow availability might not yet be reflected in the redirection of demand flows of winter tourism that further burden lower-altitude destinations while benefitting higher-altitude destinations. 

\subsection{Climate Change and Political Attitudes}
Beyond the direct effects of snow depth on party vote shares, the results speak to a broader literature on climate change and political attitudes, particularly regarding the conditions under which climate vulnerability and exposure to climate shocks shapes environmental preferences and support for green policies. 
As mentioned above the analysis shows no significant effect of snow depth on Green vote shares, as well as no significant effect of the snow anomaly on both party's vote shares. Although the results show no indication of punishment for the Green Party, as they formed part of the governing coalition of the previous government, voters still may reward the FPÖ for prioritizing short-term economic relief for vulnerable voters. As the distributive conflict and economic grievances channel far-right support, environmental salience is actively suppressed by this dynamic: material insecurity in climate-vulnerable municipalities does not translate into green voting, but rather into demand for short-term economic relief. 
The evidence opposes the findings of \textcite{baccini2021, Pink2023} who find an increased salience of environmental preferences as well as support for the Green Party in the aftermath of climate shocks. This divergence from the literature, however, points to a theoretically meaningful distinciton between acute climate shocks and gradual, slow-moving climate stress. While destructive natural disasters produce sudden, visible distruption that voters can attribute to a political cause, the long-term decline of snowfall operates below the threshold of political salience. Therefore the findings suggest that the political consequences of climate change are contingent on the nature of the climate stressor: acute shocks may heighten environmental preferences, whereas structural decline redirects grievances toward economic relief rather than environmental concern. This is further supported by the operationalization of snow anomalies as a form of \textit{climate shock}, equivalent to a natural disaster, which fails to provide significant effect - even when treated as a form of shock, the relatively small absolute magnitude of seasonal snow anomalies and its lack of disruption remains below the threshold at which climate shocks translate into visible pro-environmental behaviour. 

Furthermore, findings in the literature show that personal experience is rather short-lived \parencite{HEALY2009, egan2012}. Consequently, residents of vulnerable municipalities might not compare contemporaneous snow conditions to the ones of previous decades, but rather evaluate the most recent winter season in isolation, limiting the extent to which long-term snow decline as a climate shock shapes their electoral behavior.
Previous research on climate change and its political consequences largely leverage natural disasters with destructive consequences and the reporting of political attitudes in the immediate aftermath \parencite{egan2012, baccini2021, Masiero2021, TaheriHosseinkhani2026}. Furthermore, mechanisms of retrospective voting are often found to be short-lived and focus on post-disaster relief efforts \parencite{HEALY2009}. In the context of this study, the national elections take place five months after the end of the winter season 23/24 and there are no short-term relief measures that could alleviate the pressure of the climate shock of decreased snowfall, making it hard for voters to express support or punishment for reconstruction. The gradual decline and shift of snowfall along the Alpine regions, requires long-term and interregional climate policies, as well as investment-heavy green transitions and sustainable tourism practices that cannot be rewarded in the aftermath of one winter season. Voters make their voting decision based on the matter that is most present and salient for them, which might rather be the distributive and economic pressure than the long-term environmental prospect that is detrimental to their income base.


\subsection{Theoretical Implications and Limitations}
As shown above, the findings of the study carry several theoretical implications for the intersecting literature on economic voting, political cleavages and consequences of climate change. The results expand the literature on distributive conflicts and economic voting by showing how climate-induced resource scarcity can drive material insecurity that shapes electoral outcomes and far-right backlash. Furthermore, it shows support for a potential environment-economy trade-off, challenging the assumption that personal exposure to climate change automatically translates into increased environmental salience and green voting. The gradual decline of snow depth operates differently than sudden climate shocks such as natural disasters and therefore entail different political consequences. 

However, the scope and design of the analysis entail certain limitations that constrain the generalizability of the presented findings: 
The aggregation of vote shares at the municipal-level could mask differing individual-level preferences that balance out when aggregated. Furthermore, the sample is limited to a single country and election cycle, affecting the generalizability of the findings beyond the institutional, climatic and economic context of Austria. Future research could yield valuable insights from comparing countries that suffer or benefit differently from changing snowfall patterns. To further study the impacts of changing climate conditions over time, leveraging panel data including changes in snow availability and multiple election cycles could prove valuable. Another limiting factor is the use of national election results. Apart from different influences that shape voting decisions in general elections, the timing of Austria's national elections in September 2024 could possibly lower the salience of climate vulnerability, as it is not in the immediate aftermath of the winter season. Additionally, vote choice in national elections reflects a wide range of political preferences and considerations beyond climate grievances, limiting the extent to which snow availability alone can be isolated as a determinant of electoral outcomes, although most confounding factors are controlled for by employing district fixed effects. Future research might therefore benefit from regional elections that take place in the immediate aftermath of a winter season to avoid the dissipation of temporary grievances. To gain further knowledge on the mechanism of environmental salience, the use of survey data or climate-policy referendum outcomes could provide a more nuanced understanding of environmental preferences and their driving forces. To further explore climatic drivers of political attitudes and the extents of economic voting in the context of climate vulnerability, individual surveys and survey experiments could provide further insights . 
Moreover, as compensatory measures, such as artifical snow production, represent a major increase in expenditures for municipalities, an avenue for future research might be to investigate the effects of climate-induced compensation costs on political preferences and how it drives mechanisms of economic voting and distributive conflicts. 

\begin{figure}[H]
  \centering
  \includegraphics[width=0.85\textwidth]{plots/greens_share_map.png}
  \caption{\textit{Greens Vote Share Across Municipalities}}
  \label{fig:greens_map}
\end{figure}
\vspace{2em}




% ── VII. CONCLUSION ────────────────────────────────────────────
\section{Conclusion}
The study set out to assess the question to what extent, if at all, snow availability in Austrian municipalities influences far-right and Green party vote shares in general elections. 
Building on previous literature, I argue that snow scarcity, especially in lower-altitude municipalities, will shape party vote shares of the far-right and Green Party in general elections. I hypothesize the effect to operate through mechanisms of economic grievances driving far-right support as opposed to environmental salience due to exposure to climate change. Based on municipal-level data on average seasonal snow depth in Austrian municipalities, I measure the effect on party vote shares in Austria's national elections. The results show a modest increase in far-right vote shares in municipalities with lower seasonal snow depth. As declining snow availability is associated with less snow days per winter season, tourism-dependent municipalities experience economic grievances and material insecurity that drive far-right support in vulnerable municipalities. At the same time, the results show no significant effect of snow scarcity on Green party vote shares, indicating that even though municipalities are immediately exposed to climate change, the consequences do not translate into heightened environmental salience in electoral behaviour. Although highly vulnerable and exposed  to climate change, municipalities voted in favor of short-term economic relief as, among others, promoted by the far-right than of environmental policies and green energy transitions. This suggests that in times of economic downturn and material insecurity, voters prioritize fiscal relief policies as opposed to policies that might address the core of the problem but yield no economic benefits. The findings speak to questions around economic voting and the economic channel of far-right backlash \parencite{Rho2017, LewisBeck2019, Habersack2024}. When disproportionately burdened with material losses and insecurity, rural and vulnerable voters tend to favour far-right parties while opposing or show less support for environmental and green policies \parencite{Halikiopoulou_Vrakopoulos_Arndt_2026}. The results therefore indicate support for a trade-off between the environment and the economy \parencite{Shum2012, Mildenberger2017}. In times of economic downturn, economic preferences are of higher salience although voters are in contexts of personal exposure and economic vulnerability to climate change. While voters have multiple, often contrasting political preferences that determine vote choice in general elections \parencite{Zaller1992}, climate-induced economic grievances show the most salience to be translated into a significant effect in far-right vote shares. 
Based on climatic research, snowfall patterns continuously shift to favour higher-altitude elevations while endangering snow availability in lower-altitude destinations \parencite{li2025, Olefs2020}. Therefore I also explore the spatial heterogeneity of the effect of snow depth on party vote shares. As the analysis shows, there is no significant effect heterogeneity with changing altitudes. Lower-altitude destinations can therefore compensate for poor snow conditions through measures such as snowmaking, but at a higher cost. Yet, this does not lead to greater dissatisfactions with political decisision compared to higher-altitude municipalities. Conversely, higher-altitude destination might not yet profit enough from changing snowfall patterns and redirected demand structures to show lower salience of economic relief and higher environmental salience. 

As climate change continues to shape the natural and economic landscape of Alpine economies, understanding its consequences in the economic, political, and cultural context becomes a necessity to effectively adapt and mitigate individual grievances, but also to form a supportive voter base for climate change mitigation policies. Therefore, one immediate question coming out of this is to what extent do these findings have implications for policy making and environmental protection, as climate change continues to threaten snow as a crucial resource base for the tourism industry that relies on its availability as well as cold temperatures. Forming support for environmental protection and green transitions is not merely a question of electoral success of parties, but a structural necessity to protect the availability of natural resources that form the backbone of Alpine economies. In order to gain support for environmental protection policies, it proves to be crucial to address and prioritize immediate economic grievances of the voter base in addition to just environmental policies. Equally important is mitigating the disproportionate costs of green policies but also consequences of  climate change that vulnerable municipalities are burdened with, as these distributive conflicts shape political cleavages and far-right backlash, driving support away from the very policies needed to address the underlying crisis of climate change. 

\newpage

% ── BIBLIOGRAPHY ─────────────────────────────────────────────
\section*{Bibliography}
\printbibliography[heading=none]


\newpage

% ── APPENDIX A: SUMMARY STATISTICS ───────────────────────────
\section*{Appendix A: Summary Statistics}
\label{app:A}

\begin{table}[H]
  \caption{\textit{Descriptive Statistics of Main Variables}}
  \label{tab:summary}
  \begin{table}
\centering
\begin{tblr}[         %% tabularray outer open
]                     %% tabularray outer close
{                     %% tabularray inner open
colspec={Q[]Q[]Q[]Q[]Q[]Q[]},
hline{2}={1-6}{solid, black, 0.05em},
hline{1}={1-6}{solid, black, 0.08em},
hline{12}={1-6}{solid, black, 0.08em},
column{1}={}{halign=l},
column{2-6}={}{halign=r},
}                     %% tabularray inner close
& Mean & SD & Min & Max & N \\
fpö\_share & \num{31.97} & \num{6.91} & \num{0.00} & \num{66.66} & 2115 \\
greens\_share & \num{5.78} & \num{2.87} & \num{0.00} & \num{22.11} & 2115 \\
anomaly & \num{-0.01} & \num{0.10} & \num{-0.30} & \num{0.60} & 2115 \\
mean\_2324 & \num{0.17} & \num{0.32} & \num{0.00} & \num{2.06} & 2115 \\
altitude & \num{696.92} & \num{488.50} & \num{113.01} & \num{2643.36} & 2115 \\
log\_pop & \num{7.62} & \num{0.97} & \num{3.76} & \num{12.62} & 2115 \\
pop\_over65 & \num{21.27} & \num{3.82} & \num{11.00} & \num{40.50} & 2115 \\
pop\_foreign & \num{10.14} & \num{7.12} & \num{0.20} & \num{69.50} & 2115 \\
unemp\_rate & \num{3.69} & \num{2.24} & \num{0.00} & \num{21.60} & 2115 \\
share\_tertiary\_edu & \num{11.65} & \num{5.14} & \num{2.30} & \num{47.00} & 2115 \\
\end{tblr}
\end{table}

\end{table}

\newpage

% ── APPENDIX B: ROBUSTNESS ────────────────────────────────────
\section*{Appendix B: Robustness}
\label{app:B}

\begin{table}[H]
  \caption{\textit{Robustness: Maximum Snow Depth (2023/24)}}
  \label{tab:robustness}
  \begin{table}
\centering
\begin{talltblr}[         %% tabularray outer open
entry=none,label=none,
note{}={+ p \num{< 0.1}, * p \num{< 0.05}, ** p \num{< 0.01}, *** p \num{< 0.001}},
]                     %% tabularray outer close
{                     %% tabularray inner open
colspec={Q[]Q[]Q[]Q[]Q[]},
hline{2}={1-5}{solid, black, 0.05em},
hline{20}={1-5}{solid, black, 0.05em},
hline{1}={1-5}{solid, black, 0.08em},
hline{23}={1-5}{solid, black, 0.08em},
column{2-5}={}{halign=c},
column{1}={}{halign=l},
}                     %% tabularray inner close
& OLS FPÖ & FE FPÖ & OLS Greens & FE Greens \\
(Intercept) & \num{58.470}*** &  & \num{9.821}*** &  \\
& (\num{4.294}) &  & (\num{1.294}) &  \\
max\_2324 & \num{-1.511}*** & \num{-0.551} & \num{0.192}** & \num{0.033} \\
& (\num{0.297}) & (\num{0.351}) & (\num{0.063}) & (\num{0.074}) \\
altitude & \num{0.003}*** & \num{0.001} & \num{-0.001}*** & \num{-0.001}*** \\
& (\num{0.001}) & (\num{0.001}) & (\num{0.000}) & (\num{0.000}) \\
turnout & \num{-0.257}*** & \num{-0.257}*** & \num{-0.083}*** & \num{-0.038}** \\
& (\num{0.040}) & (\num{0.045}) & (\num{0.012}) & (\num{0.014}) \\
log\_pop & \num{1.146}*** & \num{0.479}+ & \num{0.162}** & \num{0.217}** \\
& (\num{0.203}) & (\num{0.258}) & (\num{0.060}) & (\num{0.076}) \\
pop\_over65 & \num{-0.203}*** & \num{-0.456}*** & \num{-0.103}*** & \num{0.030}* \\
& (\num{0.038}) & (\num{0.054}) & (\num{0.011}) & (\num{0.013}) \\
pop\_foreign & \num{-0.181}*** & \num{-0.131}** & \num{0.019}+ & \num{0.007} \\
& (\num{0.031}) & (\num{0.039}) & (\num{0.011}) & (\num{0.009}) \\
unemp\_rate & \num{-0.248}** & \num{-0.171} & \num{-0.187}*** & \num{-0.021} \\
& (\num{0.082}) & (\num{0.107}) & (\num{0.023}) & (\num{0.026}) \\
share\_tertiary\_edu & \num{-0.724}*** & \num{-0.662}*** & \num{0.372}*** & \num{0.397}*** \\
& (\num{0.025}) & (\num{0.045}) & (\num{0.011}) & (\num{0.014}) \\
Num.Obs. & \num{2115} & \num{2078} & \num{2115} & \num{2078} \\
R2 & \num{0.352} & \num{0.596} & \num{0.620} & \num{0.767} \\
R2 Adj. & \num{0.349} & \num{0.578} & \num{0.618} & \num{0.757} \\
\end{talltblr}
\end{table}

\end{table}
\newpage


% ── APPENDIX C: Supplementary Material ────────────────────────────────────

\section*{Appendix C: Supplementary Material}

\noindent All raw data files, as well as replication materials for the analysis can be found here: 

\noindent https://github.com/franziskarogg/snow\_voting

\end{document}